\subsection{Enlace LoRaWAN}
\label{subsec:enlace}

Esta seção descreve a configuração do enlace de rádio, o formato das mensagens trocadas
nos dois fluxos, o preparo e a fragmentação da imagem e o tratamento dos quadros
recebidos pelo servidor. Os mesmos parâmetros de rádio foram usados nos dois fluxos.

\subsubsection{Estudo preliminar e escolha da configuração de rádio}
\label{subsubsec:estudo_preliminar_enlace}

Na fase inicial do trabalho, a viabilidade do envio de imagens foi estimada
analiticamente para o plano AU915 com canais de 125~kHz e \gls{sf}~8. Nessa configuração,
o Ato n\textordmasculine~14.448 da ANATEL enquadra o \gls{lora} como sistema de salto em
frequência, que deve usar no mínimo 35 canais e limitar o tempo médio de ocupação de cada
canal a 0,4~s em um intervalo de 14~s \cite{anatel2017}. Para respeitar esse limite, cada
pacote foi restrito a 115~bytes na camada física; com isso, uma imagem JPEG de 3.182~bytes
exigiria 38 pacotes de 348,7~ms cada, totalizando cerca de 13,3~s de rádio ativo.

Essa configuração foi abandonada na implementação por dois motivos. Primeiro, o salto
entre pelo menos 35 canais é incompatível com o receptor disponível, que escuta um único
canal (Seção~\ref{subsec:simplificacoes}). Segundo, o limite de 400~ms por transmissão
restringe o tamanho do bloco de dados: um pacote com 181~bytes de \textit{payload},
tamanho adotado adiante, ocuparia 543~ms de rádio com 125~kHz.

Por isso, o enlace foi configurado para operar com largura de banda de 500~kHz, no
\textit{data rate} DR6 do plano AU915, que corresponde a SF8 com 500~kHz
(Tabela~\ref{tab:uplink_datarate}) e a um \textit{payload} máximo de 222~bytes
\cite{radioenge2022}. Com essa largura de banda, o rádio se enquadra no item~10.3 do mesmo
Ato, destinado a sistemas de modulação digital (Seção~\ref{subsec:restricoes_regulatorias}),
cujos três requisitos são verificados a seguir: largura de faixa a 6~dB de, no mínimo,
500~kHz; potência de pico máxima de 1~W; e densidade espectral de potência máxima de 8~dBm
em qualquer faixa de 3~kHz.

O módulo Radioenge RD49C transmite com potência máxima de +20~dBm (100~mW), sem comando AT
que permita ajustá-la \cite{radioenge2022}, abaixo do limite de 1~W. Supondo que essa potência se distribua de forma plana pelos 500~kHz, a densidade
espectral fica em torno de
$20 - 10\log_{10}(500/3) \approx -2{,}2$~dBm em 3~kHz, abaixo do limite de 8~dBm. Por ser
enquadrado no item~10.3, o sistema dispensa o salto de frequência e o limite de ocupação
por canal exigidos para os canais de 125~kHz, o que permite usar um canal fixo. O módulo é
homologado pela ANATEL (certificado n\textordmasculine~02021-18-07215), mas duas ressalvas
se aplicam: a largura de faixa a 6~dB do sinal \gls{lora} com 500~kHz nominais não consta
do manual e deve ser confirmada nos ensaios de homologação, e antenas com ganho superior a
6~dBi exigiriam a redução da potência (item 10.5).

O canal escolhido foi o de 915,9~MHz, o primeiro canal de 500~kHz do plano AU915: canal
64 na numeração da especificação \cite{lora_alliance_rp002_105}, que começa em zero, e
canal 65 na numeração do manual do módulo, que começa em um. No módulo, ele foi
selecionado pela máscara de canais (\texttt{AT+CHMASK}), e o \textit{data rate} foi fixado
em DR6 com o ajuste adaptativo de taxa (ADR) desativado, pois o comando \texttt{AT+DR} só
tem efeito com o ADR desligado \cite{radioenge2022}. A Tabela~\ref{tab:parametros_enlace} reúne os parâmetros
usados por todos os componentes do enlace.

\begin{table}[H]
    \centering
    \caption{Parâmetros do enlace LoRaWAN}
    \label{tab:parametros_enlace}
    \begin{tabularx}{\textwidth}{@{}lX@{}}
        \toprule
        \textbf{Parâmetro} & \textbf{Valor} \\
        \midrule
        Plano de canais & AU915 \\
        Canal & Fixo, 915,9~MHz (canal 64), sem salto de frequência \\
        \textit{Data rate} & DR6: \gls{sf}~8, \gls{bw} de 500~kHz, ADR desativado \\
        Camada física & \gls{cr} 4/5, preâmbulo de 8 símbolos, CRC ativado, cabeçalho
        explícito, palavra de sincronismo \texttt{0x34} \\
        Potência de transmissão & +20~dBm (fixa no módulo) \\
        Ativação & \gls{abp} \\
        \textit{Payload} máximo em DR6 & 222~bytes \\
        Portas de aplicação & FPort~1 (blocos de imagem) e FPort~2 (leitura do \gls{ocr}) \\
        Bloco de imagem & 175~bytes de dados + 6~bytes de cabeçalho = 181~bytes \\
        Intervalo entre pacotes & 7~s \\
        \bottomrule
    \end{tabularx}
    \legend{Fonte: Elaborado pelos autores (2026).}
\end{table}

\subsubsection{Formato das mensagens}
\label{subsubsec:formato_mensagens}

Cada mensagem \gls{lorawan} acrescenta 13~bytes ao conteúdo da aplicação: o cabeçalho MAC
(1~byte), o endereço do dispositivo (4~bytes), o campo de controle (1~byte), o contador de
quadros FCnt (2~bytes), a porta de aplicação FPort (1~byte) e o código de integridade MIC
(4~bytes). O conteúdo da aplicação (FRMPayload) é criptografado com a chave de sessão de
aplicação (AppSKey). Os dois fluxos usam portas diferentes, o que permite ao servidor
distinguir as mensagens sem nenhum outro campo. A Tabela~\ref{tab:formato_payload}
apresenta o formato do conteúdo de cada porta.

\begin{table}[H]
    \centering
    \caption{Formato do \textit{payload} de aplicação em cada fluxo}
    \label{tab:formato_payload}
    \begin{tabularx}{\textwidth}{@{}llX@{}}
        \toprule
        \textbf{Campo} & \textbf{Tamanho} & \textbf{Conteúdo} \\
        \midrule
        \multicolumn{3}{@{}l}{\textit{FPort 1 --- bloco de imagem (fluxo em nuvem)}} \\
        Identificador da leitura & 2 bytes & Número aleatório sorteado a cada captura \\
        Índice do bloco & 2 bytes & Posição do bloco na imagem, a partir de 0 \\
        Total de blocos & 2 bytes & Número de blocos da imagem \\
        Dados & até 175 bytes & Trecho do arquivo JPEG \\
        \midrule
        \multicolumn{3}{@{}l}{\textit{FPort 2 --- leitura (fluxo em borda)}} \\
        Identificador da leitura & 2 bytes & Número aleatório sorteado a cada leitura \\
        Indicadores & 1 byte & Bit 7: janela dos dígitos encontrada; bits 0 a 6: número
        $N$ de dígitos \\
        Dígitos & $N$ bytes & Um caractere ASCII por rolete: \texttt{0} a \texttt{9};
        \texttt{?} para dígito não classificado; \texttt{-} quando a janela não foi
        encontrada \\
        \bottomrule
    \end{tabularx}
    \legend{Fonte: Elaborado pelos autores (2026).}
\end{table}

Os campos numéricos são gravados em \textit{big-endian}. O identificador da leitura
impede que blocos de capturas diferentes sejam misturados na remontagem, e o par
índice--total permite ao servidor saber quais blocos faltam. No fluxo em borda, com o
medidor de seis roletes usado nos testes, a mensagem tem $3 + 6 = 9$~bytes e é enviada
em um único \textit{uplink}.

\subsubsection{Preparo da imagem e fragmentação}
\label{subsubsec:preparo_imagem}

No fluxo em nuvem, a câmera é configurada na resolução QQVGA ($160 \times 120$ pixels) em
tons de cinza, de modo que cada quadro bruto ocupa 19.200~bytes (um byte por pixel). Esse
quadro é comprimido em JPEG pelo próprio \gls{esp32cam}, por software, com qualidade 90
em uma escala de 1 a 100.

A escolha do formato foi apoiada por uma prova de conceito em C++ com a biblioteca
OpenCV, que mediu o tamanho de três imagens de hidrômetro em $160 \times 120$ pixels
após a compressão JPEG em diferentes qualidades, em três versões: colorida, em tons de
cinza e binarizada pelo método de Otsu. A Tabela~\ref{tab:poc_compressao} mostra as
médias obtidas para algumas qualidades.

\begin{table}[H]
    \centering
    \caption{Tamanho médio do JPEG de três imagens de $160 \times 120$ pixels}
    \label{tab:poc_compressao}
    \begin{tabular}{@{}cccc@{}}
        \toprule
        \textbf{Qualidade JPEG} & \textbf{Colorida (bytes)} & \textbf{Tons de cinza (bytes)}
        & \textbf{Binarizada (bytes)} \\
        \midrule
        30 & 2.028 & 1.626 & 2.773 \\
        50 & 2.338 & 1.949 & 3.433 \\
        80 & 2.910 & 2.518 & 4.725 \\
        90 & 3.466 & 3.073 & 5.895 \\
        \bottomrule
    \end{tabular}
    \legend{Fonte: Elaborado pelos autores (2026).}
\end{table}

A conversão para tons de cinza reduziu o arquivo em relação à imagem colorida, enquanto a
binarização o aumentou em quase o dobro, pois as bordas abruptas entre preto e branco são
pouco compressíveis pelo JPEG. Por isso, a imagem é transmitida em tons de cinza, sem
binarização. A qualidade 90 foi mantida para que a compressão não prejudicasse o
reconhecimento dos dígitos no servidor. Nas 60 leituras das duas corridas de medição do fluxo em nuvem, o JPEG variou de 3.323
a 4.568~bytes (mediana de 4.323~bytes), pois o tamanho depende da cena e da iluminação.
Para o dimensionamento, adotou-se como referência uma imagem de 4.300~bytes, valor
definido a partir das primeiras capturas, antes das corridas de medição, e próximo da
mediana observada depois.

O arquivo JPEG é dividido em blocos de 175~bytes, que recebem o cabeçalho de 6~bytes da
Tabela~\ref{tab:formato_payload}, totalizando 181~bytes de \textit{payload}, abaixo do
máximo de 222~bytes permitido em DR6. Uma versão anterior usava blocos de 100~bytes;
o aumento para 175~bytes reduziu o número de pacotes por imagem a pouco mais da metade,
mantendo uma folga de 41~bytes em relação ao máximo.

Os blocos são enviados ao módulo de rádio, um a um, pelo comando \texttt{AT+SENDB} na
porta~1, com intervalo de 7~s entre eles. O item~10.3 não impõe intervalo mínimo entre
transmissões em um canal fixo de 500~kHz; o valor de 7~s é conservador e foi mantido para
não sobrecarregar o módulo. Se o módulo não confirmar a transmissão em até 2~s, o bloco é
reenviado, até cinco tentativas.

\subsubsection{Tempo no ar}
\label{subsubsec:tempo_no_ar}

O tempo no ar de cada pacote foi calculado pela fórmula da Semtech para o modem \gls{lora}
\cite{semtech2013an1200_13}. A duração de um símbolo é dada pela
Equação~\ref{eq:tsym_enlace}:

\begin{equation}
    \label{eq:tsym_enlace}
    T_{\text{sym}} = \frac{2^{SF}}{BW} = \frac{2^{8}}{500\,000~\text{Hz}} = 0{,}512~\text{ms}
\end{equation}

O número de símbolos do \textit{payload} é dado pela Equação~\ref{eq:nsym_enlace}:

\begin{equation}
    \label{eq:nsym_enlace}
    n_{\text{payload}} = 8 + \max\left(\left\lceil
    \frac{8 \cdot PL - 4 \cdot SF + 28 + 16 \cdot CRC - 20 \cdot H}{4 \cdot (SF - 2 \cdot DE)}
    \right\rceil \cdot (CR + 4),\ 0\right)
\end{equation}

\noindent em que $PL$ é o tamanho do pacote na camada física, em bytes; $CRC = 1$ indica
CRC ativado; $H = 0$ indica cabeçalho explícito; $DE = 0$ indica a otimização para baixas
taxas desativada; e $CR = 1$ corresponde ao código 4/5. O tempo no ar do pacote é dado pela
Equação~\ref{eq:tpacote_enlace}, com preâmbulo de $n_{\text{pre}} = 8$ símbolos:

\begin{equation}
    \label{eq:tpacote_enlace}
    T_{\text{pacote}} = (n_{\text{pre}} + 4{,}25 + n_{\text{payload}}) \cdot T_{\text{sym}}
\end{equation}

Para um bloco de imagem completo, $PL = 13 + 181 = 194$~bytes, o que resulta em
$n_{\text{payload}} = 8 + \lceil 1564/32 \rceil \cdot 5 = 253$ símbolos e
$T_{\text{pacote}} = (8 + 4{,}25 + 253) \cdot 0{,}512 \approx 135{,}8$~ms. Para a
mensagem de leitura do fluxo em borda, $PL = 13 + 9 = 22$~bytes, o que resulta em
$T_{\text{pacote}} \approx 25{,}7$~ms. A Tabela~\ref{tab:tempo_no_ar} apresenta o número
de pacotes e os tempos para os tamanhos mínimo, mediano e máximo de JPEG medidos e para a
mensagem do fluxo em borda.

\begin{table}[H]
    \centering
    \caption{Pacotes, tempo de rádio e tempo total de envio por leitura}
    \label{tab:tempo_no_ar}
    \begin{tabular}{@{}lcccc@{}}
        \toprule
        \textbf{Conteúdo} & \textbf{Tamanho (bytes)} & \textbf{Pacotes} &
        \textbf{Rádio (s)} & \textbf{Envio total (s)} \\
        \midrule
        JPEG, menor medido & 3.323 & 19 & 2,58 & 126 \\
        JPEG, mediana medida & 4.323 & 25 & 3,36 & 168 \\
        JPEG, maior medido & 4.568 & 27 & 3,57 & 182 \\
        Leitura do fluxo em borda & 9 & 1 & 0,03 & 0,03 \\
        \bottomrule
    \end{tabular}
    \legend{Fonte: Elaborado pelos autores (2026). O envio total soma o tempo no ar ao
    intervalo de 7~s entre pacotes.}
\end{table}

Os valores mostram que, no fluxo em nuvem, o tempo de rádio é pequeno diante do tempo
total de envio, que é dominado pelos intervalos de 7~s entre os pacotes. O volume
transmitido por leitura também difere em mais de duas ordens de grandeza entre os
fluxos: cerca de 4~KB de imagem contra 9~bytes de leitura. Essa diferença é a base das
hipóteses H2 e H3.

\subsubsection{Recepção, descriptografia e remontagem}
\label{subsubsec:recepcao}

O receptor LoRa32 foi configurado com os mesmos parâmetros de camada física do
transmissor (Tabela~\ref{tab:parametros_enlace}). Para cada quadro recebido com CRC
válido, ele escreve na porta serial uma linha com o conteúdo do quadro em hexadecimal e o
RSSI medido. Como não há servidor de rede, o processamento \gls{lorawan} é feito no
computador pelo \textit{script} \texttt{reconstrutor.py}, em quatro etapas:

\begin{enumerate}
    \item \textbf{Separação dos campos}: o quadro é dividido em endereço do dispositivo,
    contador FCnt, porta e conteúdo; quadros de outros endereços são descartados.
    \item \textbf{Descriptografia}: o conteúdo é descriptografado com a AppSKey, usando o
    AES-128 em modo contador definido pela especificação \gls{lorawan}, que combina o
    endereço do dispositivo e o FCnt na geração do fluxo de chaves. O MIC não é verificado,
    pois o computador não dispõe da chave de rede.
    \item \textbf{Tratamento por porta}: na porta~2, a leitura é decodificada e validada
    (tamanho coerente com o número de dígitos e apenas caracteres permitidos). Na porta~1,
    os blocos são agrupados pelo identificador da leitura, e blocos repetidos são
    ignorados.
    \item \textbf{Remontagem e registro}: quando todos os blocos de uma leitura chegam, o
    arquivo JPEG é gravado e pode seguir para o \gls{ocr}. Se algum bloco faltar, a imagem
    é gravada como parcial, com zeros no lugar dos dados ausentes. Cada leitura é
    registrada em um arquivo CSV com os blocos recebidos, o RSSI e os tempos usados nas
    medições.
\end{enumerate}

No modo autônomo, usado nas medições, nenhum dos dois firmwares retransmite mensagens
perdidas: o nó envia a leitura ou a imagem uma única vez e volta a dormir. O firmware do
fluxo em nuvem aceita um comando de reenvio de blocos específicos, mas ele só foi usado em
testes manuais. Assim, no fluxo em nuvem, a perda de um único bloco inutiliza a imagem
inteira, e, no fluxo em borda, a perda do único pacote inutiliza a leitura. Esse
comportamento é considerado na análise da taxa de entrega.
